\section{Carry Complexity}

Carry complexity measures how much nonlinear normalization is required after
the bilinear diagonal projection step. While diagonal projection aggregates
digit interactions into raw degree masses, carry normalization transports
excess mass along the degree axis. We quantify this transport using carry mass,
carry depth, carry count, and carry entropy.

\paragraph{Carry mass.}
Let \(r_k\) denote the incoming carry at degree \(k\). Carry mass is the total
positive transported carry,
\[
    \operatorname{Mass}(r) = \sum_{k \ge 1} r_k.
\]
It measures the total amount of mass moved by normalization.

\paragraph{Carry count.}
Carry count is the number of carry positions with nonzero transported mass,
\[
    \operatorname{Count}(r) = |\{k \ge 1 : r_k \ne 0\}|.
\]
It records how many degree transitions participate in carry flow.

\paragraph{Carry depth.}
Carry depth is the largest carry index with nonzero carry,
\[
    \operatorname{Depth}(r) = \max\{k : r_k \ne 0\},
\]
with value \(-1\) when no carry appears. It measures how far normalization
propagates along the degree axis.

\paragraph{Carry entropy.}
For positive carry values, define
\[
    p_k = \frac{r_k}{\sum_{\ell:r_\ell>0} r_\ell}.
\]
Carry entropy is
\[
    H(r) = -\sum_{k:r_k>0} p_k \log_2 p_k.
\]
It distinguishes concentrated carry flow from carry flow spread across many
degree positions.

\paragraph{Diagonal mass concentration.}
For raw coefficients \(c_k\), diagonal mass concentration is
\[
    \operatorname{Conc}(c) = \frac{\max_k c_k}{\sum_k c_k},
\]
with denominator treated as one when the total raw mass is zero. This statistic
measures how strongly the projected bilinear mass is concentrated on a single
degree.

Together, these quantities support empirical laws connecting diagonal mass
profiles to carry flow. A concentrated raw diagonal distribution can create
large local excess over the base threshold, while diffuse mass may distribute
normalization cost across several smaller carry events.
